% =====================================================================
%  Apart Research - Submission Template (LaTeX port of Submission_Template.docx)
%  AI Incident Response Sprint, September 2026
%
%  Compila con: pdflatex submission.tex  (dos veces, por hyperref)
%  Formato original del .docx: US Letter, margenes de 1", Arial 11 pt
% =====================================================================
\documentclass[11pt,letterpaper]{article}

% ---------- Codificacion e idioma ----------
\usepackage[utf8]{inputenc}
\usepackage[T1]{fontenc}
\usepackage[english]{babel}   % cambia a [spanish] si escriben el informe en espanol

% ---------- Pagina y tipografia (Arial -> Helvetica) ----------
\usepackage[margin=1in]{geometry}
\usepackage{helvet}
\renewcommand{\familydefault}{\sfdefault}
\usepackage{microtype}
\linespread{1.08}
\setlength{\parskip}{0.55em}
\setlength{\parindent}{0pt}

% ---------- Utilidades ----------
\usepackage{graphicx}
\usepackage{booktabs}
\usepackage{enumitem}
\usepackage{caption}
\usepackage{xcolor}
\usepackage[most]{tcolorbox}
\usepackage{titlesec}
\usepackage[hidelinks]{hyperref}
\hypersetup{colorlinks=true,linkcolor=black,citecolor=black,urlcolor=black}

\captionsetup{font=small,labelfont=bf}
\setlist[itemize]{topsep=0.2em,itemsep=0.15em,leftmargin=1.4em}
\setlist[enumerate]{topsep=0.2em,itemsep=0.15em,leftmargin=1.6em}

% Secciones en negrita con numero y punto: "1. Introduction"
\titleformat{\section}{\large\bfseries}{\thesection.}{0.5em}{}
\titleformat{\subsection}{\normalsize\bfseries}{\thesubsection}{0.5em}{}
\titlespacing*{\section}{0pt}{1.1em}{0.4em}
\titlespacing*{\subsection}{0pt}{0.8em}{0.3em}

% Enlaces subrayados, como en el .docx
\newcommand{\ulink}[2]{\href{#1}{\underline{#2}}}

% Entrada de autor: nombre + superindice + afiliacion
\newcommand{\authorentry}[3]{%
  \begin{tabular}[t]{c}#1$^{#2}$\\[0.15em]\textit{#3}\end{tabular}}

% Texto de guia de la plantilla (borrar antes de enviar)
\newcommand{\guide}[1]{\textit{#1}}

\begin{document}
\pagestyle{plain}

% =====================================================================
%  TITULO Y AUTORES
% =====================================================================
\begin{center}
  {\LARGE\bfseries PROJECT TITLE\footnote{Research conducted at the
  \ulink{https://apartresearch.com/sprints/ai-incident-response-sprint-2026-09-11-to-2026-09-13}{AI Incident Response Sprint},
  September 2026.}}

  \vspace{1.2em}

  \begin{tabular}{ccc}
    \authorentry{Juan Camilo Melgarejo Gómez}{1}{Affiliation} &
    \authorentry{Alejandro Sánchez Poveda}{2}{Affiliation} &
    \authorentry{Sergio Montoya Ramírez}{3}{Affiliation} \\[1.6em]
    \multicolumn{3}{c}{%
      \authorentry{Daniel Díaz González}{4}{Affiliation}\hspace{4em}
      \authorentry{Roger Palomeque}{5}{Affiliation}} \\
  \end{tabular}

  \vspace{1.2em}

  \textbf{With}\\
  \textbf{Apart Research}
\end{center}

% =====================================================================
%  ABSTRACT
% =====================================================================
\begin{center}\textbf{\large Abstract}\end{center}

\begin{tcolorbox}[colback=white,colframe=black!70,boxrule=0.5pt,
                  left=8pt,right=8pt,top=6pt,bottom=6pt]
\guide{Summarize your project in 150--250 words. A strong abstract lets a
reviewer understand what you did and why it matters without reading anything
else. \underline{Make sure to cover: the problem, your approach, key results,
and the main takeaway}. Polish it last: the abstract should reflect your final
results, not your initial plan.}
\end{tcolorbox}

% =====================================================================
%  >>> BORRAR ESTE BLOQUE COMPLETO ANTES DE ENVIAR <<<
% =====================================================================
\begin{tcolorbox}[colback=black!3,colframe=black!70,boxrule=0.5pt,
                  left=8pt,right=8pt,top=6pt,bottom=6pt]
\textbf{How to use this template}: Replace the \textit{italicized guidance text}
under each section with your content. The section structure is strong guidance
but not rigid. If your project requires a different organization, feel free to
adapt. \underline{Delete all guidance text including this info box before
submitting}. Make sure the project title and author information above are up to
date.

\textbf{How your submission will be evaluated}: Your project will be judged on
the quality of this written report. To understand what we look for, see our
\ulink{https://apartresearch.notion.site/sprint-evaluation-rubric}{Evaluation Rubric}.

\textbf{Recommended length}: 4 pages excluding references and appendix. Rough
guide: Intro \& Related Work 1p, Methods and Results 2.5p, Discussion 0.5p.
\end{tcolorbox}
% =====================================================================
%  >>> FIN DEL BLOQUE A BORRAR <<<
% =====================================================================

% =====================================================================
\section{Introduction}

\guide{What problem are you addressing and why does it matter?}

\guide{Connect it to your work: we want to know why your work is practically
valuable.}

\guide{Provide enough background for readers to understand your work.}

\guide{If relevant, briefly describe the threat model or failure mode you're
addressing; reference prior work that motivates why it is worth addressing, or
explain it yourself.}

\guide{Aspire to clearly list your most important contributions that go beyond
what exists today.}

\textbf{\textit{Our main contributions are:}}
\begin{enumerate}
  \item \guide{[First contribution --- what new thing did you create, discover,
        or demonstrate?]}
  \item \guide{[Second contribution]}
  \item \guide{[Third contribution, if applicable]}
\end{enumerate}

 Early warnings of unwanted behaviour are a crucial component of the stack for controlling, tracking, and containing agents. Moreover, the cost and efficiency of these kinds of tools become a crucial aspect when considering the increasing use of Multi-Agent Systems (MAS), where unintended behaviour can originate not only within each agent individually, but also from the interactions within the system. Furthermore, the OpenAI-Hugging Face incident demonstrates that channel detection could be an unfruitful route because a swarm of agents can use increasingly diverse forms of communication, such as directory names, URL paths, or webhooks, among others, to convey messages \cite{openai2026incident}. In this context, this document proposes a methodology for detecting agent communication in MAS while circumventing the system-agnostic channel detection challenge. The principal strategy is avoiding channel detection itself and rather considering agent- and system-level entropy for detecting disruptions in the expected entropy behaviour of the MAS. This tool is not expected to replace current state-of-the-art alignment investigations, but rather to complement them with an early signal of unintended behaviour, triggering an escalation in the surveillance of behaviour for a set of agents. Our main contributions are:

\begin{enumerate}
\item \textbf{Adaptation of the MAS entropy framework for disruption detection:} By leveraging the recent definition and adaptation of entropy to agent behaviour detection \cite{zhao2026doesmultiagentcollaborationhelp, shannon_hallucination} into the general field of disruption behaviour, we present a new and computationally cheap tool for detecting unintended behaviour. This constructs an early warning sign that could trigger more advanced but costly tools, such as chain-of-thought monitoring and other state-of-the-art tooling.
\item \textbf{Development of an early benchmark for the tool:} We create a controlled experimental setting to evaluate the usefulness of the tool for disruption detection. Concretely, the benchmark simulates a situation of necessary communication and coordination between agents to complete a task by constructing a scenario of mutually asymmetric information.
\end{enumerate}

% =====================================================================
\section{Related Work}

\guide{What prior work is most similar, and how does your work differ?}

\guide{Cite the most relevant papers, tools, or projects. Explain what gap your
work addresses.}

\guide{Some questions which may help:}
\begin{enumerate}[label=(\alph*)]
  \item \guide{When and why would someone use your method over the existing
        state-of-the-art?}
  \item \guide{What information/insight does your method provide which we did
        not have before?}
\end{enumerate}

% =====================================================================
\section{Methods}

\guide{Describe your approach clearly enough that someone could replicate it.
Include key design choices and justify them where relevant (hint: the more you
can back up your design choices by referencing prior work, the better).}
The design of the experiment was reduced to the simplist case to induce the event of interest: measuring a jump or noticable change in entropy exactly when the agents are given access to communicate with each other. In the context of measuring entropy, we needed to measure the case when the agents are given access to more information, maintaining the null hypothesis that agents exposed to more information would not change significantly the entropy and the alternative hypothesis is if the quality of that information is more relevant to the task at hand, there would be a noticable jump in the change of entropy.

\guide{What models/datasets/tools did you use and why? What were key parameters
or design decisions? What did you try that didn't work? Could someone reproduce
your work from this description?}
We originally scoped this to use open models that we can run locally on our computers thinking that scoping it down to 3 simple tasks would be something easily achievable but we noticed that qwen3-8b 4 bit quantized out of the box with thinking enabled blew through our original token budget limits for accomplishing the task. We then pivoted to using OpenRouter for access to models that expose logprobs so we can use it within our experiments. So one key parameter here is having access to a model intelligent enough to use tool calling out of the box without having the harness do work nudging it along with exposing logprobs so that we can actually measure entropy.  

% =====================================================================
\section{Results}

\guide{Present your main findings with appropriate evidence. Use figures and
tables where appropriate (we strongly encourage at least one figure, see tips
below). Distinguish between observations and interpretations.}

\guide{Argue why your claims are robust. E.g., if your approach ``performs
better'' than alternatives:}
\begin{enumerate}[label=(\alph*)]
  \item \guide{Do you have enough data? Is the difference statistically
        significant?}
  \item \guide{Is it robust? Or do small changes to your setup cause substantial
        changes to your results?}
\end{enumerate}

% ---- Ejemplo de figura: descomenta y reemplaza el archivo ----
% \begin{figure}[htbp]
%   \centering
%   \includegraphics[width=0.85\linewidth]{figures/figure1.pdf}
%   \caption{Descriptive caption that can be understood without the main text.}
%   \label{fig:main}
% \end{figure}

% ---- Ejemplo de tabla ----
% \begin{table}[htbp]
%   \centering
%   \caption{Descriptive caption.}
%   \label{tab:main}
%   \begin{tabular}{lrr}
%     \toprule
%     Condition & Metric A & Metric B \\
%     \midrule
%     Baseline  & 0.00 & 0.00 \\
%     Ours      & 0.00 & 0.00 \\
%     \bottomrule
%   \end{tabular}
% \end{table}

\textbf{\textit{Tips for figures and tables:}}
\begin{itemize}
  \item \guide{Number all figures (Figure 1, Figure 2\ldots) and tables
        (Table 1, Table 2\ldots)}
  \item \guide{Include descriptive captions that can be understood without the
        main text}
  \item \guide{Place figures/tables near where they're first referenced}
  \item \textbf{\textit{Ensure text in figures is legible!}}
\end{itemize}

% =====================================================================
\section{Discussion and Limitations}

\guide{Discuss the broader implications for AI safety.}

\guide{What do your results mean? What trends do you notice and what might they
indicate?}

\subsection*{Limitations}

\guide{What are the limitations of your work? What threat models or failure
modes did you not address? Be honest about constraints --- methodological
limitations, scope limitations, or aspects you couldn't fully address in the
hackathon timeframe. Explicitly note the assumptions you made, whether
implicitly or explicitly, and how the interpretation of your results would
change if a given assumption did not hold.}

\subsection*{Future Work}

\guide{What are the natural next steps? How could this work be extended?}

% =====================================================================
\section{Conclusion}

\guide{Briefly summarize your main findings and their implications
(1--2 paragraphs).}

% =====================================================================
\section*{Code and Data}

\guide{Include links if applicable. If your project doesn't involve code
(e.g., policy analysis) or if there are info-hazard considerations, note that
here.}

\begin{itemize}
  \item \textbf{\textit{Code repository}}: \guide{[Link to GitHub/GitLab if applicable]}
  \item \textbf{\textit{Data/Datasets}}: \guide{[Link if applicable]}
  \item \textbf{\textit{Other artifacts}} \guide{(optional): [Demo link, video
        walkthrough, Hugging Face Space, etc.]}
\end{itemize}

% =====================================================================
\section*{Author Contributions (optional)}

\guide{[e.g., ``J.C.M. led the project and designed experiments. A.S.P.
implemented the code. All authors contributed to writing and reviewed the final
manuscript.'']}

% =====================================================================
\section*{References}

\guide{Use a consistent citation format. Include: Author(s), Year, Title,
Venue/Publisher, and URL or DOI where available.}

% ---- Alternativa con BibTeX: borra el enumerate de arriba y usa ----
\bibliographystyle{plain}
\bibliography{references}

% =====================================================================
\section*{Appendix \textit{(optional)}}

\guide{Supplementary material such as additional figures, detailed methodology,
prompts used, extended results, etc.}

% =====================================================================
\section*{LLM Usage Statement}

\guide{If you used LLM assistance in developing your project or writing this
report, briefly note how. Ensure all claims and results have been verified.}

\guide{NOTE: We strongly encourage that the final version of the submission is
primarily written by your team.}

\guide{[e.g., ``We used Claude to brainstorm approaches and help draft sections.
All results and claims were independently verified.'']}

\end{document}
